\chapter{Evaluation}
\label{chap:evaluation}

This chapter presents the results and evaluates them. Section~\ref{sec:tier}
tests the architectural premise, Section~\ref{sec:lambda-analysis} explains the
mechanism behind the failure, Section~\ref{sec:decoupling} evaluates the
correction, and Section~\ref{sec:qualitative} assesses the work as a whole
including what it does not establish.

Every result reports the floor a majority-class classifier achieves. Where the
evaluation protocol imposes one, the structural ceiling is reported too.

\section{Does the team tier contribute?}
\label{sec:tier}

PerMFL's premise is that it applies when known team structures exist across
devices. The direct test is to vary the team structure and measure the effect.

Two experiments were run. The first compares team assignments: teams matching
the true attack domains, teams assigned at random, and teams derived from client
updates during training. The second removes the tier by setting the team count
to one, holding the total local computation constant so that the comparison is
not confounded by training volume.

Across three datasets the result is the same. On CICIDS2017, moving from one
team to ten changes personalised macro F1 by 0.0002. On EMNIST-10, the base
paper's own dataset at its own parameter setting, the change is 0.0000. On
NSL-KDD it is 0.0011. Against a measured noise floor of 0.0116, all three are
unmeasurable.

Assignment makes no difference either. Teams matching the true day-domains do
not beat random assignment on any dataset, and derived teams do not beat either.
The largest difference across five such comparisons is 0.0016.

One configuration does show an effect. At $\lambda = 1.5$ the tier contributes
0.0342 macro F1 on CICIDS2017, above the threshold. But the paper's own
convergence condition requires $\gamma > 2\lambda$, which at $\gamma = 1.5$
caps $\lambda$ below 0.75. The setting where the tier does something lies
outside the region the theory permits. Within that region the contribution is
at most 0.0221, and operating there costs 0.157 macro F1 relative to the tuned
setting where the tier does nothing.

\section{Why: one parameter, two jobs}
\label{sec:lambda-analysis}

The explanation is in the update rule rather than in the results.

$\lambda$ appears twice. In the device update \eqref{eq:device} it scales the
penalty pulling a device toward its team, so a smaller value permits more
personalisation. In the team update \eqref{eq:team} it weights $\bar{\theta}_i$,
the average of the team's members, which is the only term in that equation
carrying any information about who is in the team.

The relative weight of that term is $\lambda\eta / \eta\gamma = \lambda/\gamma$.
At the tuned setting of $\lambda = 0.05$ with $\gamma = 1.5$, the members'
average is weighted 0.0015 against a pull of 0.045 toward the global model, so
membership influences the team model by 3.3 per cent. The team model is
effectively a copy of the global model, which is why the number of teams makes
no difference.

The convergence condition bounds this. Since $\gamma > 2\lambda$ implies
$\lambda/\gamma < 0.5$, a team model can never take more than half its update
from its own members within the region the theory permits.

Tracing the consequence for the global model: substituting $w = x$ at the start
of a round gives $w \approx 0.9985x + 0.0015\bar{\theta}$, and after the global
update the server model moves by roughly 0.00135 of the direction its devices
indicate. Over ten rounds it travels about 1.35 per cent from its
initialisation. The measured global accuracy of 0.0495 on nine classes, below
the 0.111 a uniform random classifier achieves, is consistent with a model that
has not left its starting point.

The two roles want opposite values. Personalisation wants $\lambda$ small; a
functional global model wants it comparable to $\gamma$. The shipped
parameterisation forces a choice.

\section{Decoupling the two roles}
\label{sec:decoupling}

Separating them is a one-parameter change. The device penalty keeps $\lambda$;
the team update takes an independent weight.

\subsection{Headline comparison}

On CICIDS2017 with attack-domain partitioning, benign traffic thinned to
twenty-five per cent of the training split, and evaluation on unseen classes,
across five paired seeds:

\begin{center}
\begin{tabular}{lrrrr}
\toprule
metric & shipped & decoupled & difference & paired $t$ \\
\midrule
Personal macro F1 & $0.2809 \pm 0.0017$ & $0.3154 \pm 0.0044$ & $+0.0345$ & 21.28 \\
Personal accuracy & $0.7842 \pm 0.0062$ & $0.8161 \pm 0.0042$ & $+0.0320$ & 10.61 \\
Global macro F1 & $0.1139 \pm 0.0345$ & $0.2155 \pm 0.0187$ & $+0.1016$ & 4.35 \\
Global accuracy & $0.5412 \pm 0.1885$ & $0.8487 \pm 0.0041$ & $+0.3075$ & 3.60 \\
\bottomrule
\end{tabular}
\end{center}

All four exceed the critical value of 2.776 and win on every seed. The floors
are 0.8171 accuracy and 0.0999 macro F1. The shipped global model sits below the
accuracy floor; the decoupled one sits above it.

\subsection{Variance}

The most consistent effect is not the mean. Across-seed standard deviation in
global accuracy falls from 0.1885 to 0.0041 on CICIDS2017, from 0.0928 to 0.0004
on TON-IoT under host partitioning, and from 0.0619 to 0.0063 on NSL-KDD. Two to
three orders of magnitude, and it holds in the one configuration where the mean
regressed.

The shipped global model varies between roughly 0.35 and 0.73 depending on
initialisation. For a security product that is arguably the more consequential
finding: an unpredictable detector cannot be operated against a threshold.

\subsection{Across heterogeneity}

Sweeping client heterogeneity from near-homogeneous to strongly skewed, three
seeds per point:

\begin{center}
\begin{tabular}{lrrrrrr}
\toprule
configuration & JSD & ceiling & $\Delta$PM & $t$ & $\Delta$GM & $t$ \\
\midrule
CICIDS $k{=}1$ & 0.1730 & 0.308 & $+0.0273$ & 5.45 & $+0.1209$ & 3.49 \\
CICIDS $k{=}3$ & 0.1414 & 0.596 & $+0.1160$ & 32.49 & $+0.3973$ & 44.74 \\
CICIDS $k{=}5$ & 0.0851 & 0.795 & $+0.2306$ & 36.39 & $+0.5130$ & 57.91 \\
CICIDS $k{=}8$ & 0.0001 & 1.000 & $+0.3256$ & 82.83 & $+0.7049$ & 213.4 \\
TON-IoT $\alpha{=}0.1$ & 0.7591 & 0.781 & $+0.0391$ & 5.84 & $+0.1572$ & 7.78 \\
TON-IoT $\alpha{=}0.5$ & 0.3909 & 0.985 & $+0.1093$ & 17.44 & $+0.2881$ & 8.86 \\
TON-IoT $\alpha{=}5.0$ & 0.0646 & 1.000 & $+0.1614$ & 20.01 & $+0.3764$ & 26.76 \\
NSL-KDD $\alpha{=}0.1$ & 0.5379 & 1.000 & $+0.1365$ & 12.98 & $+0.1592$ & 5.53 \\
NSL-KDD $\alpha{=}5.0$ & 0.0333 & 1.000 & $+0.1181$ & 16.29 & $+0.1128$ & 5.75 \\
\bottomrule
\end{tabular}
\end{center}

The correction improves both models at every level on all three datasets, and
the improvement grows as heterogeneity falls. On CICIDS2017 the personalised
gain rises from 0.027 at the most skewed setting to 0.326 at the least, and the
global gain from 0.121 to 0.705.

That direction was not what the project expected, and the earlier expectation is
worth recording as having been wrong. The reasoning had been that the team
weight matters more when members diverge. The measurements say the opposite, and
the explanation is straightforward in hindsight: the term the weight controls is
the average of the team's members, and an average of divergent members carries
little usable signal. When members agree, the same term carries a great deal.
The correction lets the model use information it already had.

At the least skewed CICIDS setting the decoupled configuration reaches 0.8432
personalised and 0.8546 global macro F1 against a ceiling of 1.0, compared with
0.5176 and 0.1497 for the shipped configuration.

\subsection{A regression}

One configuration goes the other way. Under attack-domain partitioning on
TON-IoT, at the highest heterogeneity measured anywhere in the study, the global
macro F1 falls from 0.0917 to 0.0403, $t = -6.21$, losing on all five seeds. The
decoupled global accuracy is 0.2369 with a standard deviation of zero, which is
exactly the majority-class rate: the model collapsed to predicting normal
traffic.

This bounds the claim. Making the global model listen harder to strongly
specialised teams can destroy it.

\section{Local training as a baseline}
\label{sec:local}

A baseline with no federation at all, each client training alone on its own
data, outperforms every federated configuration on personalised macro F1 at
every heterogeneity level tested. On CICIDS2017 it scores 0.9504 against the best
federated result of 0.9379.

Two qualifications matter. When clients are evaluated on classes absent from
their training data, the local baseline loses more than half its score, from
0.9525 to 0.4012 at one setting, because a class absent from the loss cannot be
predicted. And the local baseline produces no global model at all, so it offers
nothing to a host with no local data or one that has not yet seen a given attack.

The honest conclusion is that federated personalisation does not pay on this
data, and that the value of the method lies in its global model. That makes the
finding in Section~\ref{sec:lambda-analysis} more consequential rather than
less: the shipped parameterisation cripples precisely the component that
justifies federating.

\section{Critical evaluation}
\label{sec:qualitative}

\subsection{What the work establishes}

The reproduction is exact and the defect catalogue is verifiable line by line,
so neither depends on trusting the analysis. The tier finding is supported
across three datasets and four team counts with a mechanism derived from the
update rule rather than inferred from outcomes, and the mechanism predicts the
regression in Section~\ref{sec:decoupling} as well as the improvements. The
decoupling result is paired, seeded and significant across twelve
configurations.

\subsection{What it does not establish}

The comparison against other methods is absent. Neither SCMoE-PFL nor CFMD-i
publishes an implementation, so no relative claim about accuracy or efficiency
is made. Runs completed in minutes on one CPU core where a mixture-of-experts
implementation required hours on a GPU, but the client counts, datasets and
software stacks differ, so that observation is structural rather than measured.

The baselines shipped with the base implementation were not run on the intrusion
detection data. Only the method under test carries the metric instrumentation
this study requires, and extending it to the others was not completed. The
consequence is that the results position PerMFL against its own configurations
and against local training, not against the federated literature.

Several exploratory results rest on single runs. The local depth and
communication sweeps are reported as exploratory for that reason, though their
effects, up to 0.29 macro F1, are far outside the noise floor.

The clustering signal is parameter-space only. CFMD-i additionally uses an
output-space divergence, which is likely the better signal under label skew,
since two clients can hold near-identical weights and behave differently. Given
that the term the clustering feeds is weighted at 3.3 per cent, a better signal
would probably not change the conclusion, but that is an argument rather than a
measurement.

\subsection{Threats to validity}

The heterogeneity measure is the divergence between label distributions, which
does not capture feature-space or concept differences. A partition producing
identical label distributions with different decision boundaries would register
as homogeneous.

The structural ceiling calculation assumes a client can only predict classes it
trained on. A model that generalises across classes could exceed it, so the
bound is conservative rather than exact.

The per-client cap flattens client sizes to a narrow band, suppressing quantity
skew that real deployments exhibit. The host-based partition is the exception,
with client sizes spanning 1,984 to 37,832 rows.

Two datasets share an origin, and CICIDS2017 and NSL-KDD are both laboratory
captures with scripted attacks. Agreement between them is weaker evidence than
agreement between independent sources would be.
